
\section{Phân tích và xác thực Input}

\subsection{Tổng quan về xử lý input}

Chương trình cần xác thực input từ người dùng ở hai bối cảnh chính:
\begin{enumerate}
    \item \textbf{Giai đoạn Setup}: Đọc 4 tọa độ khi đặt tàu.
    \item \textbf{Giai đoạn Battle}: Đọc 2 tọa độ khi tấn công.
\end{enumerate}

Xác thực input đảm bảo chương trình không bị lỗi khi người dùng nhập sai định dạng hoặc giá trị ngoài phạm vi.

\subsection{Thủ tục parse\_four\_ints}

Thủ tục này tách 4 số nguyên từ một chuỗi input. Nó được sử dụng khi đặt tàu.

\subsection{Mục tiêu}

Chuyển đổi một chuỗi nhập vào (ví dụ: ``0 0 0 1'') thành 4 giá trị số nguyên.

\subsection{Thuật toán}

\begin{verbatim}
PROCEDURE parse_four_ints(input_string, addr_val1, addr_val2, 
                          addr_val3, addr_val4)
  parsed_count = 0
  current_number = 0
  current_char_index = 0
  is_building_number = false
  
  WHILE current_char_index < length of input_string:
    char = input_string[current_char_index]
    
    IF char is a digit (0-9):
      // Add digit to current number
      digit_value = char - '0'
      current_number = current_number * 10 + digit_value
      is_building_number = true
    
    ELSE IF char is whitespace (space, tab, newline):
      IF is_building_number:
        // Number is complete, store it
        SWITCH parsed_count:
          CASE 0: *addr_val1 = current_number
          CASE 1: *addr_val2 = current_number
          CASE 2: *addr_val3 = current_number
          CASE 3: *addr_val4 = current_number
        
        parsed_count += 1
        current_number = 0
        is_building_number = false
        
        IF parsed_count == 4:
          RETURN success
    
    ELSE:
      // Invalid character (not digit, not whitespace)
      // In some implementations, this could trigger an error
      // For simplicity, we skip invalid characters
      ;
    
    current_char_index += 1
  
  // Handle last number if exists
  IF is_building_number AND parsed_count < 4:
    SWITCH parsed_count:
      CASE 0: *addr_val1 = current_number
      CASE 1: *addr_val2 = current_number
      CASE 2: *addr_val3 = current_number
      CASE 3: *addr_val4 = current_number
    parsed_count += 1
  
  IF parsed_count == 4:
    RETURN success
  ELSE:
    RETURN failure

END PROCEDURE
\end{verbatim}

\subsection{Chi tiết hoạt động}

Thủ tục này hoạt động theo các bước sau:

\begin{enumerate}
    \item \textbf{Khởi tạo}: Đặt parsed\_count = 0, current\_number = 0, và chỉ số vị trí ký tự = 0.
    
    \item \textbf{Quét từng ký tự}: Lặp qua từng ký tự trong chuỗi input.
    \begin{itemize}
        \item Nếu ký tự là chữ số: Nhân current\_number với 10 và cộng giá trị chữ số mới (để xây dựng số multi-digit).
        \item Nếu ký tự là khoảng trắng: Nếu đang xây dựng số, lưu số đó vào địa chỉ tương ứng, tăng parsed\_count, và đặt lại current\_number.
        \item Nếu ký tự khác: Bỏ qua (hoặc có thể báo lỗi).
    \end{itemize}
    
    \item \textbf{Xử lý số cuối cùng}: Sau vòng lặp, nếu vẫn đang xây dựng số, lưu nó.
    
    \item \textbf{Kiểm tra}: Nếu parsed\_count == 4, trả về thành công. Ngược lại, thất bại.
\end{enumerate}

\subsection{Ví dụ}

Input: ``2 3 2 4''

\begin{enumerate}
    \item Quét '2': digit, current\_number = 2.
    \item Quét ' ': whitespace, lưu 2 vào addr\_val1, parsed\_count = 1, current\_number = 0.
    \item Quét '3': digit, current\_number = 3.
    \item Quét ' ': whitespace, lưu 3 vào addr\_val2, parsed\_count = 2, current\_number = 0.
    \item Quét '2': digit, current\_number = 2.
    \item Quét ' ': whitespace, lưu 2 vào addr\_val3, parsed\_count = 3, current\_number = 0.
    \item Quét '4': digit, current\_number = 4.
    \item End of string: lưu 4 vào addr\_val4, parsed\_count = 4.
    \item Trả về success.
\end{enumerate}

\subsection{Thủ tục parse\_two\_ints}

Thủ tục này tương tự parse\_four\_ints nhưng chỉ tách 2 số nguyên. Được sử dụng khi tấn công.

\begin{verbatim}
PROCEDURE parse_two_ints(input_string, addr_val1, addr_val2)
  // Tương tự parse_four_ints nhưng lặp đến khi parsed_count == 2
  
  parsed_count = 0
  current_number = 0
  current_char_index = 0
  is_building_number = false
  
  WHILE current_char_index < length of input_string:
    char = input_string[current_char_index]
    
    IF char is digit:
      digit_value = char - '0'
      current_number = current_number * 10 + digit_value
      is_building_number = true
    
    ELSE IF char is whitespace:
      IF is_building_number:
        SWITCH parsed_count:
          CASE 0: *addr_val1 = current_number
          CASE 1: *addr_val2 = current_number
        
        parsed_count += 1
        current_number = 0
        is_building_number = false
        
        IF parsed_count == 2:
          RETURN success
    
    current_char_index += 1
  
  IF is_building_number AND parsed_count < 2:
    SWITCH parsed_count:
      CASE 0: *addr_val1 = current_number
      CASE 1: *addr_val2 = current_number
    parsed_count += 1
  
  IF parsed_count == 2:
    RETURN success
  ELSE:
    RETURN failure

END PROCEDURE
\end{verbatim}

\subsection{Xác thực phạm vi (Range Validation)}

Sau khi tách số nguyên từ chuỗi input, chương trình phải kiểm tra các giá trị có nằm trong phạm vi cho phép hay không.

\subsection{Khi đặt tàu}

\begin{verbatim}
// Kiểm tra 4 tọa độ có nằm trong [0, 6]
IF (r1 < 0 OR r1 > 6 OR c1 < 0 OR c1 > 6 OR
    r2 < 0 OR r2 > 6 OR c2 < 0 OR c2 > 6):
  PRINT "Error: Coordinates must be in range [0, 6]!"
  RETURN failure
\end{verbatim}

\subsection{Khi tấn công}

\begin{verbatim}
// Kiểm tra 2 tọa độ có nằm trong [0, 6]
IF (row < 0 OR row > 6 OR col < 0 OR col > 6):
  PRINT "Error: Coordinates must be in range [0, 6]!"
  RETURN failure
\end{verbatim}

\subsection{Xử lý input sai}

Khi input không hợp lệ, chương trình:

\begin{enumerate}
    \item \textbf{Thông báo lỗi rõ ràng}: In thông báo lỗi cụ thể (ví dụ: ``Invalid format'', ``Out of range'', v.v.).
    
    \item \textbf{Yêu cầu nhập lại}: Quay lại bước đọc input, yêu cầu người dùng nhập lại.
    
    \item \textbf{Không cập nhật trạng thái}: Không thay đổi bất kỳ dữ liệu nào trên bàn cờ hoặc bản đồ bắn.
\end{enumerate}

\subsection{Vòng lặp input-validate-retry}

\begin{verbatim}
LOOP:
  PRINT "Enter coordinates: "
  READ input_line
  
  CALL parse_two_ints(input_line, addr_row, addr_col)
  
  IF parse failed:
    PRINT "Invalid input format. Please try again."
    CONTINUE LOOP
  
  IF row < 0 OR row > 6 OR col < 0 OR col > 6:
    PRINT "Coordinates out of range [0, 6]. Please try again."
    CONTINUE LOOP
  
  // Valid input, proceed
  BREAK LOOP
\end{verbatim}

\subsection{Giới hạn hiện tại}

Chương trình xác thực input có những giới hạn sau:

\begin{itemize}
    \item \textbf{Không xử lý số âm}: Input có số âm sẽ không được phân tích đúng. Ví dụ, ``-1 2'' sẽ không được tách thành -1 và 2.
    \item \textbf{Không xử lý số quá lớn}: Nếu input chứa số lớn hơn 32-bit integer, hành vi không xác định.
    \item \textbf{Không xử lý ký tự đặc biệt}: Input chứa ký tự đặc biệt (trừ khoảng trắng) sẽ bị bỏ qua hoặc gây lỗi.
\end{itemize}

Những giới hạn này là chấp nhận được vì input được giả định là từ người chơi hợp lý, không phải từ input tương tác phức tạp.

\subsection{Best Practice}

Khi viết code MIPS để xác thực input:

\begin{enumerate}
    \item \textbf{Tách biệt parse và validate}: Parse (tách số) riêng, validate (kiểm tra phạm vi) riêng.
    \item \textbf{Cung cấp feedback rõ ràng}: Thông báo lỗi phải chỉ ra chính xác vấn đề gì.
    \item \textbf{Lặp lại cho đến hợp lệ}: Không bỏ qua input sai; yêu cầu người dùng nhập lại.
    \item \textbf{Kiểm thử kỹ lưỡng}: Test với input hợp lệ, không hợp lệ, và biên (boundary).
\end{enumerate}
